\documentclass[11pt]{article}
\usepackage[margin=1in]{geometry}
\usepackage{acronym}
\usepackage{booktabs}

\title{Acronym Expansion Stress Case}
\author{Test Corpus}
\date{}

\begin{document}
\maketitle

\begin{abstract}
This fixture tests first-use acronym expansion with the \texttt{acronym} package
and combines the expansions with ordinary section and table cross-references.
\end{abstract}

\section{Acronym Register}
\label{sec:register}

\begin{acronym}[SUTVA]
\acro{ATE}{average treatment effect}
\acro{RCT}{randomized controlled trial}
\acro{SUTVA}{stable unit treatment value assumption}
\end{acronym}

\section{Identification Language}
\label{sec:identification}

The first reference to an \ac{RCT} should expand the term before showing the
short form. The first reference to the \ac{ATE} should behave similarly. Later
references to the \ac{RCT} and the \ac{ATE} should use the short forms because
the terms have already appeared in the document.

The design discussion also invokes \ac{SUTVA}. Section~\ref{sec:register}
contains the acronym list, and Table~\ref{tab:acronym-check} summarizes the
terms used in the paper.

\begin{table}[htbp]
\centering
\caption{Acronyms used in the fixture}
\label{tab:acronym-check}
\begin{tabular}{ll}
\toprule
Short form & Concept \\
\midrule
\acs{RCT} & Trial design \\
\acs{ATE} & Causal estimand \\
\acs{SUTVA} & Interference condition \\
\bottomrule
\end{tabular}
\end{table}

\section{Conclusion}
\label{sec:conclusion}

The final paragraph refers again to \ac{SUTVA}, which should now appear in short
form. This gives the staged corpus a compact test for package-managed acronym
state across sections.

\end{document}
